\section{Gaussian 混合模型让类别成为隐变量}
\label{sec:13-Machine-Learning/06-Probabilistic-Models/03}

一个客服系统记录了三万通电话的通话时长，你需要一个密度模型来判断哪些通话异常。画出直方图，形状明显有两个峰：一堆很短的查询类通话挤在两分钟附近，另一堆处理类通话散在六七分钟。用单个 Gaussian 去拟合，最大似然给出的均值落在两峰之间的低谷上——那个位置几乎没有电话，而模型认为它最典型。加大方差也救不了，只会让两个峰同时被压平。

数据里显然有两种通话，可是记录里没有``通话类型''这一列。前两节的生成式分类都假设标签已知；现在把标签遮住，模型还能不能自己把它认出来？

\subsection{把不可观测的成分编号当成一个随机变量}

设想每个样本先掷一次骰子选定一个隐藏成分 \(Z_i\)，再由该成分生成观测：

\[
P(Z_i=k)=\pi_k,\qquad
X_i\given Z_i=k\sim\mathcal N(\mu_k,\Sigma_k).
\]

把 \(Z_i\) 边缘掉，剩下的就是我们真正能看见的密度：

\[
p(x_i)=\sum_{k=1}^K\pi_k\,\mathcal N(x_i\given\mu_k,\Sigma_k).
\]

\begin{definition}[Gaussian 混合模型]
由混合权重 \(\pi_k\)、成分均值 \(\mu_k\) 与协方差 \(\Sigma_k\) 定义的上述潜变量模型称为 Gaussian 混合模型，其边缘密度是各成分密度的加权和。
\end{definition}

这一步的意义超出高斯本身。左边是一个形状复杂、无法用任何标准分布描述的密度；右边是几个最简单的椭球密度做加权和。隐变量在这里扮演的角色是把复杂的边缘分布拆成简单成分的叠加，主题模型对文档做同样的事，因子分析对连续隐向量做同样的事（见\hyperref[sec:13-Machine-Learning/04-Dimensionality-Reduction/06]{``因子分析与概率 PCA：把低秩变成生成模型''}）。所以 \(Z_i\) 不必对应任何真实存在的东西，它是一件建模装置，用来换取写得出来的密度。

\begin{center}
\begin{tikzpicture}[>=stealth]

% 直方图
\foreach \a/\h in {0/0.06, 0.5/0.26, 1/0.78, 1.5/1.42, 2/1.20, 2.5/0.52,
                   3/0.72, 3.5/1.02, 4/0.82, 4.5/0.32, 5/0.11, 5.5/0.03}
  \fill[black!10] (\a,0) rectangle (\a+0.5,\h);

\draw[gray!60,->] (-0.1,0) -- (6.8,0);
\draw[gray!60,->] (0,-0.1) -- (0,1.95);
\node[font=\scriptsize,anchor=north] at (6.6,-0.05) {通话时长};

% 两个成分
\draw[blue!60,densely dotted,thick]
  plot[domain=0:6.4,samples=110] (\x,{1.55*exp(max(-(\x-1.8)^2/0.72,-9))});
\draw[red!60,densely dotted,thick]
  plot[domain=0:6.4,samples=110] (\x,{1.05*exp(max(-(\x-3.8)^2/0.50,-9))});

% 混合密度
\draw[violet,very thick]
  plot[domain=0:6.4,samples=140]
  (\x,{1.55*exp(max(-(\x-1.8)^2/0.72,-9))+1.05*exp(max(-(\x-3.8)^2/0.50,-9))});

% 单个 Gaussian 的最大似然拟合
\draw[black,dashed,thick]
  plot[domain=0:6.4,samples=110] (\x,{0.92*exp(max(-(\x-2.75)^2/3.2,-9))});

\node[font=\scriptsize,violet] at (1.15,1.82) {混合密度};
\node[font=\scriptsize] at (5.45,1.05) {单个 Gaussian};
\draw[black!45] (5.45,0.95) -- (4.6,0.55);
\end{tikzpicture}
\end{center}

\emph{单峰模型把均值放在谁也不在的地方；两个简单成分相加就把形状接住了。}

图里那条虚线正是开头的困境：一个成分的模型只能在两个峰之间折中。紫色曲线由两条点线相加而成，每一条都只是普通的一维 Gaussian——复杂度不在成分里，在``几个成分怎样加起来''这件事里。

\subsection{软归属来自 Bayes 公式，但它不构成学习}

参数给定时，样本 \(i\) 来自成分 \(k\) 的后验概率是

\[
\gamma_{ik}=P(Z_i=k\given x_i)
=\frac{\pi_k\,\mathcal N(x_i\given\mu_k,\Sigma_k)}
{\sum_j\pi_j\,\mathcal N(x_i\given\mu_j,\Sigma_j)}.
\]

\(\gamma_{ik}\) 通常叫 responsibility，读作成分 \(k\) 对这个样本负多少责任。它和 K-means 的硬指派形成对照：落在两峰之间那些通话时长为四分钟的电话，会得到 \(0.5\) 上下的责任值，模型如实承认自己分不清。K-means 必须把它交给某一个簇，代价是丢掉这份不确定性，细节见\hyperref[sec:13-Machine-Learning/05-Clustering-and-Data-Geometry/01]{``K-means 在交替下降中制造簇''}。

麻烦在于这个公式要求参数已知。反过来，若每个 \(Z_i\) 已知，参数也很好办：按成分分组，各自算加权均值和加权协方差，与有标签的 GDA 完全一样。两个方向各自简单，可现在两端都空着。直接写出观测的 log-likelihood

\[
\ell(\theta)=\sum_{i=1}^n\log\Bigl[\sum_{k=1}^K\pi_k\,\mathcal N(x_i\given\mu_k,\Sigma_k)\Bigr]
\]

就能看出症结：对数套在求和外面，没法像有标签时那样按成分拆开，逐个成分的更新式也就写不出来。下一节的 EM 正是从这个卡点出发，利用两端各自容易的结构来回交替。

\subsection{混合模型的解天生不唯一，而且可以坏到无界}

即使 EM 收敛得很好，也要先弄清楚它收敛到的是什么。把所有成分的编号同时置换，\(\pi\) 与 \((\mu_k,\Sigma_k)\) 跟着一起换，混合密度一个字都不变。似然函数因此有 \(K!\) 个完全等价的最大值点，``第 3 个成分''这种说法在跨运行时没有意义：今天的第 3 号可能就是昨天的第 1 号。要比较两次拟合，得先按参数距离做匹配，再谈对齐后的差异。这是标签置换不可辨识，它不是数值问题，无法靠更好的初始化或更长的迭代消除。

更危险的是似然可以无界。取某个成分，把 \(\mu_k\) 放到某个具体样本点 \(x_i\) 上，再让 \(\abs{\Sigma_k}\to0\)。该点处的 Gaussian 密度与 \(\abs{\Sigma_k}^{-1/2}\) 成正比，趋于无穷；其余成分照常解释其他样本，于是整体 log-likelihood 也趋于无穷。这类退化解在数据里有孤立点时特别容易被找到：一个成分塌缩到那个孤立点上，最小特征值一路压到你设的下界，那个点的责任值锁死在 \(1\)。把下界继续调小，训练似然还能继续涨，而模型什么也没学到。所以``训练似然更高''在混合模型里不是一个可信的进步信号。

实践中的对策是给协方差设地板、共享或对角化协方差、加一个 Inverse-Wishart 类的先验做惩罚。这些都是在把无界的搜索空间截断成有界的。此外多个随机初值不可省略：似然是非凸的，局部极小与鞍点普遍存在（关于非凸地形上迭代法能保证什么，见\hyperref[sec:09-Convex-Optimization/08-Nonconvex-and-Stochastic-Optimization/01]{``下降引理在非凸地形中只保证接近驻点''}）。反复重启后你会发现，训练似然最高的那个解与留出似然最好的那个解，经常不是同一个。

\subsection{成分数不能从训练拟合里白拿}

增大 \(K\) 几乎不会降低最大训练似然，因为小模型总能被大模型模仿。多出来的自由度会被用去追逐局部涨落，切出一些只覆盖十几个样本的细碎成分。可用的判据有留出对数似然、BIC、以及重抽样下的解稳定性，但它们回答的是不同问题：预测密度最好的 \(K\)、描述长度最短的 \(K\)、分群最稳定的 \(K\) 常常不相等，彼此相差一两个成分是常态。

诊断时值得盯住几个信号：有没有权重 \(\pi_k\) 接近零的成分，有没有协方差的最小特征值贴着地板，多个初值给出的是相近的密度还是仅仅相同的标签编号。还有一个便宜的检验：用已知参数生成合成数据，看拟合能否在允许标签置换后把参数找回来；如果在自己造的数据上都恢复不了，在真实数据上的成分就更没有解释力。

需要强调的是，这些诊断针对的是``模型是否被合理地拟合''，不针对``成分是否对应真实类型''。开头那批通话被分成两个成分，不等于客服业务里就存在两种通话；换一个特征表示，或者多给一个成分，划分方式可能完全不同。混合模型给出的是一种能把密度接住的分解，它的可解释性需要在模型之外获得证据。带着这层保留，我们回到被搁置的那个卡点：对数里的求和该怎么处理。
